% TOP_PROBLEM: {"id": 2631, "title": "Kourovka Notebook Problem 21.122", "category_id": 17, "category": {"id": 17, "name": "group_theory", "display_name": "Group Theory", "description": "Problems about groups, group actions, representations, and related algebraic structures.", "slug": "group-theory", "order_index": 17, "created_at": "2026-07-31T15:26:25.671Z"}, "status": "open", "problem_number": "KOU-21.122", "set_id": 10}
% FIRST_POSTED: 2026-10-01
% ENTRY_KIND: statement_only
% UnsolvedMath ID 2631; source catalogue code KOU-21.122
% !TeX program = lualatex
% Definition and sources only; this file is not an LLM solution attempt.
\documentclass[11pt,a4paper]{article}
\usepackage[margin=25mm]{geometry}
\usepackage{fontspec}
\setmainfont{lmroman10-regular.otf}[BoldFont=lmroman10-bold.otf,
 ItalicFont=lmroman10-italic.otf,BoldItalicFont=lmroman10-bolditalic.otf]
\usepackage{amsmath,amssymb,mathtools}
\usepackage{unicode-math}
\setmathfont{latinmodern-math.otf}
\AtBeginDocument{\let\setminus\smallsetminus}
\usepackage{xurl,hyperref}
\hypersetup{colorlinks=true,urlcolor=blue}
\setcounter{secnumdepth}{0}
\setlength{\parindent}{0pt}
\setlength{\parskip}{5pt}
\setlength{\emergencystretch}{3em}
\begin{document}
\section*{Kourovka Notebook Problem 21.122}\label{2631}
{\small UnsolvedMath ID 2631 · KOU-21.122 · Group Theory · Kourovka Notebook - New Problems, Issue 21}
\subsection{Definitions and mathematical statement}
A profinite group is a compact Hausdorff topological group expressible as an inverse limit of finite discrete groups. It need not be metrizable or topologically finitely generated. A word $w=w(x_1,\ldots,x_r)$ is a fixed element of the free group on finitely many symbols; substitution defines a continuous word map
\[
 w:G^r\longrightarrow G,
 \qquad (g_1,\ldots,g_r)\longmapsto w(g_1,\ldots,g_r).
\]
Let $G_w=w(G^r)$ denote its image, the set of word values. Since $G^r$ is compact and $G$ is Hausdorff, this image is compact.

Write $\mathfrak c=2^{\aleph_0}=|\mathbb R|$ for the cardinality of the continuum. The question is whether every word and every profinite group satisfy the alternative
\[
 |G_w|<\infty\qquad\text{or}\qquad |G_w|\ge\mathfrak c.
\]
Equivalently, must $|G_w|<\mathfrak c$ imply that $G_w$ is finite?

The count concerns distinct values themselves, not the subgroup generated by them and not its topological closure. The comparison with $\mathfrak c$ is a cardinal comparison; it includes ruling out any infinite cardinal strictly below the continuum, without assuming the continuum hypothesis. The profinite group is unrestricted in size, and the word is allowed to be arbitrary, including words whose values are always the identity. Finiteness is asserted separately for the particular word and group under consideration.
\subsection{Short English statement}
Must every infinite word-value set in a profinite group contain at least continuum many elements?
\subsection{Sources}
\begin{itemize}
\item[S1] ULAM.AI and the UnsolvedMath Contributors, supplied problems.json, record 2631 (KOU-21.122), Kourovka Notebook - New Problems, Issue 21. Catalogue identity and source transcription; CC BY 4.0.
\url{https://huggingface.co/datasets/ulamai/UnsolvedMath}.
\item[S2] The Kourovka Notebook, 21st edition, version 47 (30 September 2026), new problems, Problem 21.122. Expanded formulation of the supplied catalogue statement.
\url{https://arxiv.org/pdf/1401.0300v47}.
\end{itemize}
\end{document}
